\documentclass[../main.tex]{subfiles}

\begin{document}

\subsection{Overview}

CVE-2024-45020 is a \textbf{verifier bug} in the eBPF stack-state comparison logic used during state pruning. The bug is located in \textbf{\texttt{stacksafe()}}, where the verifier compares the stack state of a newly explored path with a previously visited state to decide whether the new path can be \textbf{safely pruned}.

The vulnerability affects several Linux 6.x stable branches and is caused by an incorrect assumption during stack comparison. When the verifier compares an old state with a current state, it iterates over the stack depth allocated by the old state, but may access stack metadata that is outside the stack state allocated by the current state. This causes an \textbf{invalid memory access inside the verifier} and can lead to a verifier crash.

\begin{table}[H]
\centering
\begin{tabularx}{\textwidth}{|l|X|}
\hline
\textbf{CVE ID} & CVE-2024-45020 \\
\hline
\textbf{Affected versions} & Linux 6.6.15--6.6.47 and 6.7--6.10.6; affected 6.11 pre-release kernels before the upstream fix \\
\hline
\textbf{Component} & eBPF verifier, \texttt{stacksafe()} state pruning \\
\hline
\textbf{Vulnerability type} & Invalid memory access in \texttt{stacksafe()} leading to false pruning in the tested exploit \\
\hline
\textbf{Root cause} & Missing bound check between \texttt{old->allocated\_stack} and \texttt{cur->allocated\_stack} \\
\hline
\textbf{Impact} & Direct-map pointer and kernel heap-layout disclosure; potential exploit building block, but no complete kernel-text KASLR bypass \\
\hline
\textbf{CVSS} & 5.5 MEDIUM \\
\hline
\end{tabularx}
\caption{Summary of CVE-2024-45020.}
\end{table}

The main result of our analysis is that the vulnerability can be driven beyond the official crash condition in a BPF-capable setting. The local exploit checks for \texttt{CAP\_BPF} and documents execution as root or with BPF-related capabilities. This distinction is important: even when the process has a non-zero UID, it is not a fully unprivileged BPF loader.

Without a BPF-capable execution context, the same approach does not reliably reach the vulnerable code path. The verifier applies stricter checks to unprivileged BPF programs, including more conservative scalar precision tracking. As a result, \texttt{regsafe()} rejects the candidate states before the vulnerable stack comparison in \texttt{stacksafe()} can be reached.

\subsection{Vulnerability Analysis}

The eBPF verifier explores all possible execution paths of a program. To avoid state explosion, it uses state pruning: when execution reaches an instruction already analyzed before, the verifier compares the current state with the previously stored state. If the current state is considered equivalent or safer, the verifier skips the new path.

This comparison is performed in stages. First, register states are compared through \texttt{regsafe()}. Then, stack states are compared through \texttt{stacksafe()}.

The vulnerable logic can be summarized as follows:

\begin{lstlisting}[language=C, caption={Simplified vulnerable stacksafe logic.}]
static bool stacksafe(struct bpf_verifier_env *env,
                      struct bpf_func_state *old,
                      struct bpf_func_state *cur)
{
    for (int i = 0; i < old->allocated_stack; i++) {
        spi = i / BPF_REG_SIZE;

        if (old->stack[spi].slot_type[i % BPF_REG_SIZE] !=
            cur->stack[spi].slot_type[i % BPF_REG_SIZE])
            return false;
    }

    return true;
}
\end{lstlisting}

The bug is that the loop is bounded by \texttt{old->allocated\_stack}, but the comparison also reads \texttt{cur->stack[spi].slot\_type} \textbf{without first checking whether \texttt{cur} has allocated that stack depth}.

The vulnerable condition is:

\begin{itemize}
    \item \textbf{Path A} is analyzed first and writes deep into the stack, for example at \texttt{fp[-512]}.
    \item \textbf{Path B} is analyzed later and writes only near the top of the stack, for example at \texttt{fp[-8]}.
    \item At a merge point, \texttt{old->allocated\_stack = 512}, while \texttt{cur->allocated\_stack = 8}.
    \item \texttt{stacksafe()} iterates up to the old stack depth and reads beyond the allocated stack array of the current state.
\end{itemize}

\begin{lstlisting}[caption={Stack-depth mismatch causing the bug.}]
Path A:
  write fp[-512]
  old->allocated_stack = 512

Path B:
  write fp[-8]
  cur->allocated_stack = 8

stacksafe(old, cur):
  iterates up to 512
  reads cur->stack[1..63] out of bounds
\end{lstlisting}

Without memory debugging instrumentation, this invalid access may not immediately stop verification in the tested setup. In our exploit, the verifier-side out-of-bounds access can make the stack comparison appear successful, causing the verifier to incorrectly classify the current state as equivalent to the old one. This produces the second stage of the bug: false pruning.

\subsection{Exploit Environment}

The disclosure stage depends on residual data remaining in stack memory. The tested leak configuration therefore uses \texttt{CONFIG\_INIT\_STACK\_NONE=y}. If \texttt{CONFIG\_INIT\_STACK\_ALL\_ZERO=y} is enabled, the false-pruning condition can still be reached, but the runtime values read from the uninitialized stack are zeroed and the pointer-disclosure primitive is neutralized. The kernel configuration and the reported result must therefore be considered together.

The exploit was executed in a BPF-capable context. This is a privileged condition even when the userspace UID is not zero, and the result should not be interpreted as a fully unprivileged local information disclosure.

\subsection{Exploit Primitive}

The primitive observed in our exploit is not a direct arbitrary read/write from BPF runtime. Instead, the exploit abuses a verifier-side invalid memory access to influence pruning. If the out-of-bounds stack metadata comparison behaves as if the current state were covered by the old one, \texttt{stacksafe()} returns true and the verifier skips a path that should have been analyzed.

The effect is:

\begin{lstlisting}[caption={False pruning effect.}]
Verifier:
  believes Path B is equivalent to Path A
  skips analysis of deep stack read

Runtime:
  executes Path B
  reads from stack depth never initialized on that path
  leaks residual kernel data
\end{lstlisting}

The runtime leak occurs because the verifier approved a read from a stack location that is valid only in the old path. On the current path, the same stack depth was never initialized. Therefore, the program can copy residual data into a BPF map and return it to userspace.

The resulting primitive is mainly an \textbf{information leak}. In the successful privileged version, the exploit leaks addresses from the x86-64 kernel direct-map region. These values disclose the location and allocation pattern of kernel heap objects, making the CVE relevant as a building block for more complex exploitation chains. However, the observed values do not directly disclose the randomized kernel-text base and the primitive does not provide an arbitrary write.

\subsection{Exploit Architecture}

The exploit constructs \textbf{two asymmetric verifier paths}:

\begin{enumerate}
    \item a big-stack path that touches deep stack slots;
    \item a small-stack path that allocates only shallow stack slots.
\end{enumerate}

The two paths are later merged at a pruning point. The exploit tries to force the verifier into calling \texttt{stacksafe(old, cur)} with \textbf{incompatible stack allocation sizes}.

The privileged exploit uses a bounded loop or back-edge to create an efficient merge point. This is important because the verifier naturally revisits the same instruction with different states, allowing the pruning logic to compare the big-stack state against the small-stack state.

A simplified model is:

\begin{lstlisting}[caption={High-level exploit model.}]
1. Create Path A:
     write deep stack slot fp[-512]

2. Create Path B:
     write shallow stack slot fp[-8]

3. Force both paths to merge.

4. Trigger stacksafe(old, cur):
     old stack depth = 512
     cur stack depth = 8

5. Obtain false pruning.

6. Runtime reads uninitialized deep stack bytes.

7. Copy leaked data into a BPF map.
\end{lstlisting}

The exploit uses a BPF array map as an exfiltration channel. Once the program is accepted, runtime execution stores leaked values into the map. Userspace then reads the map and prints the leaked values.

\subsection{Experimental Results}

The execution of the target exploit produced the following log output:

\begin{figure}[H]
    \centering
    \includegraphics[width=\textwidth]{img/CVE-2024-45020.png}
    \caption{Exploit execution log for CVE-2024-45020.}
    \label{fig:log-2024-45020}
\end{figure}

The privileged PoC successfully confirms the vulnerability in the tested environment. With the right branch and pruning structure, the verifier accepts the program and runtime execution leaks values such as:

\begin{lstlisting}[caption={Examples of addresses leaked from the uninitialized BPF stack.}]
0xffff888002498ee0
0xffff888002498c20
0xffff888002498120
0xffff8880024983e0
\end{lstlisting}

On x86-64, addresses beginning with \texttt{0xffff8880} belong to the kernel \textbf{direct map}, also called the \textbf{physmap}. This virtual region maps physical RAM into the kernel address space. The leaked values are therefore pointers to memory-backed kernel objects or buffers rather than pointers to executable kernel code.

The observed addresses have a regular structure: they end at the same offset within their containing objects and many consecutive values differ by \texttt{0x2c0} bytes. This strongly suggests that the leak repeatedly exposes the same field from homogeneous objects allocated by the SLUB allocator. Since each trigger sends a datagram through a socket filter, networking objects such as \texttt{sk\_buff}-related allocations are plausible sources. Nevertheless, the address values alone are insufficient to identify the exact object type, so this attribution remains an inference rather than a confirmed result.

These leaks reveal useful information about the kernel heap:

\begin{itemize}
    \item concrete virtual addresses of kernel allocations;
    \item relative placement and spacing of objects from the same allocation pattern;
    \item evidence that heap grooming produced a predictable layout;
    \item a potential address oracle for a second UAF, OOB-write or type-confusion primitive.
\end{itemize}

They are not kernel-text pointers, which on this architecture would normally have a form such as \texttt{0xffffffff81...}. Consequently, the observed disclosure does not recover the kernel-text base and should not be described as a complete KASLR bypass. It is more accurately a \textbf{physmap and heap-layout disclosure}: valuable for reducing address uncertainty in a larger exploit chain, but insufficient by itself for arbitrary memory access or privilege escalation.

The vulnerability status can be summarized as:

\begin{itemize}
    \item verifier bug confirmed;
    \item false pruning confirmed;
    \item direct-map kernel pointer disclosure confirmed;
    \item repeatable kernel heap-layout information observed;
    \item kernel-text KASLR bypass not demonstrated by the observed values;
    \item LPE not completed;
    \item unprivileged exploitation found structurally blocked in our tests.
\end{itemize}

The most important negative result concerns unprivileged exploitation. Several variants were tested using forward-only branch structures, heap grooming and repeated retries. The exploit worked reliably when executed with the required capabilities, but not as a pure unprivileged user.

The root cause of this limitation is \texttt{env->bpf\_capable}. When the process is not BPF-capable, the verifier enables \textbf{stricter precision tracking}. As a consequence, the branch register becomes marked as precise, and \texttt{regsafe()} rejects the state comparison before \texttt{stacksafe()} is reached.

\begin{lstlisting}[caption={Unprivileged failure point.}]
Unprivileged mode:
  precise scalar tracking enabled

regsafe():
  detects different scalar ranges
  states are not equivalent
  pruning does not happen

stacksafe():
  vulnerable code path is never reached
\end{lstlisting}

This distinction is important for impact assessment. The vulnerability is real and exploitable under privileged BPF conditions, but it is not a straightforward unprivileged local root exploit in the tested configuration.

\subsection{Comparison with Previous eBPF Exploits}

Among the exploits analyzed in the previous reports, the closest comparison for CVE-2024-45020 is CVE-2023-2163. Both vulnerabilities are rooted in the verifier's pruning logic: the verifier tries to avoid re-analyzing execution states that appear equivalent to states already proved safe. In both cases, the security failure comes from accepting an equivalence relation that is not actually conservative enough.

In CVE-2023-2163, the verifier incorrectly prunes a branch because precision tracking does not preserve all dependencies needed for pointer-bound safety. The skipped path can still execute at runtime, producing a one-byte stack overflow. The exploit then brute-forces the corrupted byte, leaks a map pointer, upgrades the primitive to arbitrary 64-bit read/write and finally patches credentials to obtain root or escape a container. This exploit path is explicitly described in the previous report as a transition from incorrect path pruning to arbitrary kernel read/write and privilege escalation.

CVE-2024-45020 follows the same high-level idea of unsafe pruning, but the faulty comparison happens in a different layer of the verifier state. Instead of losing precision on a register dependency, the verifier compares two stack states with different allocation depths. The old state has a deeper allocated stack, while the current state is shallower. Because \texttt{stacksafe()} accesses metadata from the current stack before checking whether that depth is actually allocated, the comparison can read invalid verifier-heap data and lead to a false pruning decision.

\begin{table}[H]
\centering
\begin{tabularx}{\textwidth}{|l|X|X|}
\hline
\textbf{Aspect} & \textbf{CVE-2023-2163} & \textbf{CVE-2024-45020} \\
\hline
Verifier mistake & Incorrect path pruning due to incomplete precision tracking & Incorrect stack-state comparison during pruning \\
\hline
Faulty verifier area & Register precision propagation and path equivalence & \texttt{stacksafe()} and stack-depth equivalence \\
\hline
State mismatch & Runtime path has different pointer bounds than the pruned state & Current path has a smaller allocated stack than the previously visited state \\
\hline
First primitive & One-byte stack overflow & Verifier heap OOB read during stack comparison \\
\hline
Primitive upgrade & Brute-force low byte, leak map pointer, build arbitrary 64-bit R/W & False pruning, runtime read from uninitialized stack, leak direct-map pointers \\
\hline
Exploit result & Arbitrary kernel R/W, \texttt{cred} patching, LPE/container escape & Physmap and heap-layout disclosure in privileged setup; no kernel-text base or reliable arbitrary write \\
\hline
Privilege constraints & Requires BPF loading capabilities in the tested environment & Requires BPF-capable execution; unprivileged variants were blocked before reaching \texttt{stacksafe()} \\
\hline
\end{tabularx}
\caption{Exploit-level parallelism between CVE-2023-2163 and CVE-2024-45020.}
\end{table}

\subsection{Patch and Mitigation}

The vulnerability was fixed by adding a missing bounds check in \texttt{stacksafe()} before accessing the current state's stack array.

Conceptually, the fix is:

\begin{lstlisting}[language=C, caption={Conceptual model of the fix.}]
for (i = 0; i < old->allocated_stack; i++) {
    spi = i / BPF_REG_SIZE;

    if (i >= cur->allocated_stack)
        return false;

    if (old->stack[spi].slot_type[i % BPF_REG_SIZE] !=
        cur->stack[spi].slot_type[i % BPF_REG_SIZE])
        return false;
}
\end{lstlisting}

After the patch, if the current state has allocated less stack than the old state, the two states are immediately considered non-equivalent. This prevents both the verifier heap OOB read and the false pruning decision.

Recommended mitigations are:

\begin{itemize}
    \item upgrade to a kernel containing the fix or an appropriate stable backport;
    \item restrict access to \texttt{CAP\_BPF} and other capabilities that permit privileged BPF use;
    \item disable unprivileged BPF where possible;
    \item enable automatic stack initialization to limit the demonstrated residual-data disclosure; this hardening does not repair the verifier bug;
    \item enable KASAN in testing environments to detect verifier heap OOB reads;
    \item treat BPF-capable containers as high-risk isolation boundaries.
\end{itemize}

\subsection{Security Impact}

CVE-2024-45020 demonstrates that verifier vulnerabilities do not always need to corrupt runtime program memory directly. In this case, the first memory violation happens inside the verifier itself while comparing abstract states. The resulting false pruning then authorizes a runtime path that leaks kernel data.

The practical impact is mainly confidentiality. In our tests, the leaked qwords were direct-map pointers that disclose where kernel heap objects reside and how similarly sized objects are arranged. Such information can weaken heap randomization assumptions and act as an address oracle for a separate memory-corruption primitive. It does not reveal the contents of the pointed-to objects, because the leaked scalar cannot be dereferenced by the BPF program, and it does not directly reveal the randomized kernel-text base.

The disclosure could strengthen another exploit in the same controlled environment. The kernel configuration uses SLUB without freelist randomization or freelist hardening, making allocation patterns more predictable. If combined with a UAF, OOB write or pointer type confusion, the leaked address could help locate a target object and verify that heap grooming succeeded. The current CVE-2024-45020 payload alone cannot perform that upgrade: it reads only residual bytes from a fixed BPF-stack offset and provides neither arbitrary read nor arbitrary write.

The most relevant lesson is that pruning correctness is security-critical. The verifier is allowed to skip a path only if the new state is truly covered by a previously analyzed one. If this equivalence check reads invalid memory or compares states with incompatible stack depths, the verifier may approve a program whose runtime behavior exposes kernel memory.

In short, the verifier tried to optimize analysis by pruning an apparently equivalent path, but the equivalence decision was based on an out-of-bounds heap read. That mistake turns a performance optimization into a kernel heap-address disclosure primitive.

\end{document}
